\documentclass[sloppy,oribibl,envcountsame,envcountsect,11pt]{article}

\usepackage{amsfonts,amssymb,amsthm,amsmath}
\usepackage{hyperref}
\usepackage{mathtools}
\usepackage{comment}

%\usepackage{newpxtext} % T1, lining figures in math, osf in text
%\usepackage{textcomp} % required for special glyphs
%\usepackage[varg,bigdelims]{newpxmath}
%\linespread{1.05}
\usepackage[margin=1in]{geometry}

\linespread{1.1}

%\usepackage[margin=1in]{geometry}

\usepackage[utf8]{inputenc} 
\usepackage{amsmath,amssymb,amsthm}
\usepackage{algorithm}

%\theoremstyle{plain}

%\usepackage{algpseudocode}

%\usepackage{a4}


\usepackage{float}
\usepackage{verbatim}

\usepackage{todonotes}
\newcommand{\mnote}[3]{\todo[color=#3!40,size=\footnotesize]{\textbf{#2:} #1}}
\newcommand{\inote}[3]{\todo[color=#3!40,inline]{\textbf{#2:} #1}}
\newcommand{\pushkar}[1]{\mnote{#1}{P}{blue}}
\DeclareMathOperator {\range}{Range}
\DeclareMathOperator {\col}{Col}

\title{On Efficient Noncommutative Polynomial Factorization via Higman Linearization \thanks{A preliminary version was presented at 
  the 37th Computational Complexity Conference, CCC'22, \cite{AJ22}.}}

\author{V. Arvind\thanks{Institute of Mathematical Sciences, Chennai,
    India and Chennai Mathematical Institute, Siruseri, Kelambakkam,
    India, \texttt{email: arvind@imsc.res.in}} \and Pushkar S
  Joglekar\thanks{Vishwakarma Institute of Technology, Pune, India,
    \texttt{email: joglekar.pushkar@gmail.com}}
    \thanks{Author would like to thank SERB for the funding through the MATRICS project, File no. MTR/2018/001214}
    }

\date{}

\newtheorem{theorem}{Theorem}[section]
\newtheorem{corollary}[theorem]{Corollary}
\newtheorem{definition}[theorem]{Definition}
\newtheorem{lemma}[theorem]{Lemma}
\newtheorem{fact}[theorem]{Fact}
\newtheorem{claim}[theorem]{Claim}
\newtheorem{remark}[theorem]{Remark}
\newtheorem{observation}[theorem]{Observation}
\newtheorem{proposition}[theorem]{Proposition}
\newtheorem{conjecture}[theorem]{Conjecture}
\newtheorem{hypothesis}[theorem]{Hypothesis}
\newtheorem{exercise}[theorem]{Exercise}
%\newtheorem{problem}[theorem]{Open problem}

%\newenvironment{proof} {\noindent{\it Proof. }} {{\qed}}

\newenvironment{proofof}[1]{\noindent{\it Proof of #1. }} {{\qed}}

\newtheorem*{theorem*}{Theorem}

\newtheorem{problem}[theorem]{Problem}

\usepackage{enumitem}

 %% \newenvironment{problem}[2][]{\begin{plainproblem}[#1] #2%
 %% \begin{description}[nosep,leftmargin=1.90cm,%
 %%  style=sameline,align=right,format=\normalfont\emph,
 %%  beginpenalty=10000,midpenalty=5000]%
 %% }{%
 %%   \end{description}\end{plainproblem}%
 %% }


\def\qed{\hspace*{\fill} $\Box$\par\medskip}

\DeclareMathOperator{\diag}{\rm{diag}}

%
\usepackage{graphicx}

\usepackage{amsmath,amssymb}

\newcommand{\op}[1]{\ensuremath{\operatorname{#1}}}

\newcommand{\field}{\mathbb{F}}
\newcommand{\F}{\mathbb{F}}
\renewcommand{\angle}[1]{\langle #1 \rangle}
\newcommand{\FX}{\F\angle{X}}
\newcommand{\fx}{\F\angle{x_1,\ldots,x_n}}
\newcommand{\pd}{\partial}

%\newcommand{\pdl}{\partial_\ell}
%\newcommand{\pdr}{\partial_r}

\newcommand{\pdl}{\partial_{\textsc{l}}}
\newcommand{\pdr}{\partial_{\textsc{r}}}

\newcommand{\FP}{\mathrm{FP}}
\newcommand{\NC}{\mathrm{NC}}
\newcommand{\MA}{\mathrm{MA}}
\newcommand{\coNP}{\mathrm{coNP}}
\newcommand{\NP}{\mathrm{NP}}
\newcommand{\size}{\mathrm{size}}
\newcommand{\var}{\op{Var}}
\newcommand{\poly}{\mathrm{poly}}
\newcommand{\mon}{\op{Mon}}
\newcommand{\SOP}{\textsc{SOP}}
\newcommand{\rank}{\op{rank}}

\renewcommand{\span}{\op{span}}

\newcommand{\R}{\textsc{R}}

\newcommand{\Q}{\mathbb{Q}}

%\newcommand{\fR}{\mathfrak{R}}
\newcommand{\fR}{\FX}
%\newcommand{\fF}{\mathfrak{F}}

\newcommand{\skewf}{\mathbb{F}\newbrak{X}}
\newcommand{\fF}{\skewf}


\newcommand{\prob}[1]{\textsc{#1}}

\DeclareFontEncoding{LS1}{}{}
\DeclareFontSubstitution{LS1}{stix}{m}{n}
\DeclareSymbolFont{symbols2stix}{LS1}{stixfrak} {m} {n}
\DeclareMathSymbol{\lparenless}{\mathopen} {symbols2stix}{"32}
\DeclareMathSymbol{\rparengtr}{\mathclose}{symbols2stix}{"33}
 
%\newcommand{\newbrak}[1]{\angle{\angle{\angle{#1}}}}
\newcommand{\newbrak}[1]{{\lparenless} #1 {\rparengtr}}



%\usepackage{figure}
\usepackage{framed}

\begin{document}

\maketitle

\begin{abstract}

In this paper we study the problem of efficiently factorizing
polynomials in the free noncommutative ring
$\F\angle{x_1,x_2,\ldots,x_n}$ of polynomials in noncommuting
variables $x_1,x_2,\ldots,x_n$ over the field $\F$. We obtain
the following result:

\begin{itemize} 
\item[] Given a noncommutative algebraic branching
  program\footnote{This strengthens the main result in 
    earlier versions of this paper where the algorithm was only
    for noncommutative arithmetic formulas.} of size $s$ computing a
  noncommutative polynomial $f\in\F\angle{x_1,x_2,\ldots,x_n}$ as
  input, where $\F=\F_q$ is a finite field, we give a randomized
  algorithm that runs in time polynomial in $s, n$ and $\log_2q$ that
  computes a factorization of $f$ as a product $f=f_1f_2\cdots f_r$,
  where each $f_i$ is an irreducible polynomial that is output as a
  noncommutative algebraic branching program.
   
\item[] The algorithm works by first transforming the given algebraic
  branching program computing $f$ into a linear matrix $L$ using
  Higman's linearization of polynomials. We then factorize the linear
  matrix $L$ and recover the factorization of $f$. We use basic
  elements from Cohn's theory of free ideals rings combined with
  Ronyai's randomized polynomial-time algorithm for computing
  invariant subspaces of a collection of matrices over finite fields.
\end{itemize}

\noindent\textbf{Keywords: Noncommutative Polynomials, Arithmetic
  Circuits, Factorization, Identity testing.}
\end{abstract}

\newpage

\setcounter{page}{0}
\tableofcontents
\thispagestyle{empty}
\newpage

\input{ncf-intro.tex}
\input{ncf-prelim.tex}
\input{ncf-comnz.tex}
\input{ncf-comz.tex}

\section{Concluding Remarks}

In this paper we present a randomized polynomial-time algorithm for
the factorization of noncommutative polyomials \emph{over finite
  fields} that are input as \emph{algebraic branching programs}. The
irreducible factors are output as algebraic branching programs.

Several open questions arise from our work. We mention two of them.
The first question is the complexity of factorization over rationals
of noncommutative polynomials given as ABPs or arithmetic formulas. Our
approach involves the crucial use of Ronyai's algorithm for invariant
subspace comptutation which turns out to be a hard problem over
rationals.  We believe a different approach may be required for the
rational case.

The use of Higman linearization prevents us from generalizing this
approach to noncommutative polynomials given as arithmetic
circuits. We do not know any nontrivial complexity upper bound for the
factorization problem for noncommutative polynomials given as
arithmetic circuits.\\

\noindent\textbf{Acknowledgements.}~~We thank anonymous referees
and Partha Mukhopadhyay for asking about extension of our earlier
factorization algorithm for noncommutative formulas to algebraic
branching programs.


%\bibliographystyle{plain}

\bibliographystyle{alpha}
\bibliography{references}
\appendix

\newpage
\input{appendix.tex}


\end{document}
% end of file template.tex

















































\begin{lemma}\label{amitsur}[Amitsur's theorem]
Let $f \in \mathbb{F}\langle X \rangle$ be a non constant polynomial  of degree $d\geq 1$, where $X=\{x_1, x_2, \ldots, x_g \}$ be a set of noncommuting variables. Then there exist $d \times d$ matrices $M_1,\ldots, M_g$ such that $f(M_1, \ldots, M_g) \in \mathbb{GL}_{d}(\ mathbb{F})$. Moreover given $f$ by a noncommutative formula of size $s$, the matrix substitution $M_1, \ldots, M_g$ can be computed in randomized poly(s, d, g) time. 
\end{lemma}
\begin{definition}[Matrix blow-up]

\end{definition}

%\begin{lemma}\label{hkv20-theorem}[Proposition 2.7 in \cite{HKV20}]
%Let $f \in \mathbb{F}\langle X \rangle$ be a non constant polynomial  of degree $d\geq 1$, where $X=\{x_1, x_2, \ldots, x_g \}$ be a set of noncommuting variables. Let $M_1, \ldots, M_g$ be $d \times d$ be the scalar matrices  such that $A=f(M_1, M_2, \ldots, M_g)$ is an invertible scalar matrix. Let $g=A^{-1}f(Y_1+M_1, Y_2+ M_2, \ldots, Y_g + M_g)$, where $Y_i = (y^{i}_{j,k})_{1 \leq j, k \leq d}$ for free noncommuting variables $y^{i}_{j,k}$ (i.e. $g$ is matrix-blow up or point centered amplification of $f$ wrt $\overline{M}= (M_1, \ldots, M_g)$). Then $g(\overline{0})= I_{d}$ and $f$ is irreducible implies $g$ is an atom.  
%\end{lemma}

\begin{lemma}
Let $L \in (\mathbb{F}\langle X \rangle)^{n \times n}$ be a monic linear matrix with $L= A_0 + A_1 x_1 + \ldots A_g x_g$ where $X=\{x_1, x_2, \ldots, x_g \}$ be a set of noncommuting variables and $A_i \in Mat_{n \times n}(\mathbb{F})$ for $i=1$ to $g$. Let $M_1, \ldots, M_g$ be $d \times d$ scalar matrices  such that $R=L(M_1, M_2, \ldots, M_g)$ is an invertible matrix. Let $L'=R^{-1}L(Y_1+M_1, Y_2+ M_2, \ldots, Y_g + M_g) \in (\mathbb{F}\langle Y\rangle ^{nd \times nd})$, where $Y_i = (y^{(i)}_{j,k})_{\{j,k\}}$ for $1 \leq i \leq g, 1 \leq j, k \leq d$ for free noncommuting variables $y^{(i)}_{j,k}$'s  and $Y= \cup_{i=1}^g \textrm{var}(Y_i)$. Then $L$ is an atom iff $L'$ is an atom. Moreover, given $P,Q \in \mathbb{GL}_{nd}(\mathbb{F})$ such that
\[
PL'Q=  
\left(
\begin{array}{cc}
A' & 0 \\
D' & B'
\end{array}
\right)\]

where $A' \in (\mathbb{F}\langle Y \rangle)^{r' \times r'}$ and $B' \in (\mathbb{F}\langle Y \rangle)^{s' \times s'}$ for $r',s' >0$ and $r'+s'=nd$. There is an efficient deterministic algorithm with complexity polynomial in $n,d,g$ to compute $U,V \in \mathbb{GL}_{n}(\mathbb{F})$ such that 
\[
ULV=  
\left(
\begin{array}{cc}
A & 0 \\
D & B
\end{array}
\right)\]
where $A \in (\mathbb{F}\langle X \rangle)^{r \times r}$ and $B \in (\mathbb{F}\langle X \rangle)^{s \times s}$ for $r,s >0$ and $r+s=n$. 
 \end{lemma}


\section{Higman's linearization and Cohn's notion of stable associates}

Analogous to the notion of integral domains and unique factorization
domains in commutative ring theory, P.M.~Cohn \cite{,,} has developed
a theory for noncommutative rings based on the concept of weak
algorithms (which is a noncommutative generalization of the division
algorithm) and free ideal rings. We present the basic definitons and
notions of this theory relevant to our paper.

\subsection*{Higman's linearization process}

We explain the basic ideas involved along with an overview of our
algorithmic result. After that we will present the details.

Let $M$ be an $n\times n$ matrix over $\FX$ such that $M[n,n]=f+g\times h$.
Then, we apply the following steps:
\begin{itemize}
\item Expand $M$ to an  $(n+1)\times (n+1)$ matrix by adding a new last
  row and last column with diagonal entry $1$ and remaining new entries
  zero:
\[
\left[
\begin{array}{c|c}
M & 0 \\
\hline
0 & 1
\end{array}
\right]. 
\]
\item Then the bottom right $2\times 2$ submatrix is transformed as
  follows by elementary row and column operations

\[
\left(
\begin{array}{cc}
f+gh & 0 \\
0 & 1
\end{array}
\right)\rightarrow
\left(
\begin{array}{cc}
f+gh & g \\
0 & 1
\end{array}
\right)\rightarrow
\left(
\begin{array}{cc}
f & g \\
-h & 1
\end{array}
\right)
\]

\end{itemize}

Given a polynomial $f\in\FX$ by repeated application of the above
Higman linearization we finally obtain a matrix $L=A_0+\sum_{i=1}^n
A_i x_i$, where each $A_i, 0\le i\le n$ is an $r\times r$ over $\F$,
for some $r$ that depends on $f$. The Higman linearization implies
that there are $r\times r$ matrices $P$ and $Q$ with entries in $\FX$
such that

\begin{theorem}[Higman Linearization]{\rm\cite{Cohn2}}\label{higthm}
  Given a polynomial $f\in\FX$ there are square matrices $P$ and $Q$
 and a linear matrix $L$ such that
\begin{itemize}
\item  $P$ is upper triangular and $Q$ is lower triangular.
\item The diagonal entries of $P$ and $Q$ are all $1$'s (hence,
  $P^{-1}$ and $Q^{-1}$ both exist and are matrices with entries in
  $\FX$).  
\item   
\[
PAQ = L,
\]
where
\[
A :=
\left(
\begin{array}{c|c}
f & 0 \\
\hline
0 & I_{r-1}
\end{array}
\right). 
\]
\end{itemize}
\end{theorem}

Hrubes and Wigderson, in their paper \cite{HW14} initiating the study
of rational identity testing, show that if a noncommutative polynomial
$f$ is given by an arithmetic formula of size $s$ then the resulting
linear matrix $L$ after Higman linearization is $O(s)\times O(s)$.

We now state a fundamental result of Cohn which allows us to relate
the factorization of $f$ to the factorization of the matrix $L$.

\begin{theorem}{\rm\cite[Theorem 5.8.8]{fir-book}}\label{cohnthm}
  Let $C$ be an $n\times n$ linear matrix with entries from $\FX$
  such that $C$ is a monic and full matrix. Then $C$ is not irreducible if
  and only if either $n=0$ or $n>1$ and there are $n\times n$
  invertible scalar matrices $P$ and $Q$ such that
  \[
  PCQ =
  \left(
\begin{array}{cc}
A & 0 \\
D & B
\end{array}
\right)
\]
where $A$ is an $r\times r$ matrix and $B$ is an $s\times s$ matrix such
that $r+s=n$.
\end{theorem}

It follows from Theorem~\ref{higthm} that the polynomial $f$ is
irreducible if and only if the linear matrix $L$ is irreducible since
the matrices $A$ and $L$ are stable associates.

Furthermore, assuming $L$ is a monic matrix we can use Theorem~\ref{cohnthm} to
check if $L$, and hence $f$, is irreducible or not. Building on this
further, when $\F$ is a finite field, we can use Ronyai's randomized
polynomial-time algorithm \cite{Ron2} to find a common invariant
subspace of the matrices $A_0,A_1,\ldots,A_n$ which will give us
the above decomposition. Clearly, we have the factorization
\[
  \left(
\begin{array}{cc}
A & 0 \\
D & B
\end{array}
\right) =   \left(
\begin{array}{cc}
A & 0 \\
0 & I
\end{array}
\right)\cdot 
  \left(
\begin{array}{cc}
I & 0 \\
D & I
\end{array}
\right)\cdot
\left(
\begin{array}{cc}
I & 0 \\
0 & B
\end{array}
\right).
\]
We can see from the proof of Cohn's result \cite{Cohn2} that the
matrices
\[
\left(
\begin{array}{cc}
I & 0 \\
0 & B
\end{array}
\right)
\text{ and }
\left(
\begin{array}{cc}
A & 0 \\
0 & I
\end{array}
\right)
\]
are non-units while the remaining matrix is clearly a unit. Overall,
this is a nontrivial factorization of $L$. Our main result is an
algorithm to efficiently recover from this factorization a
factorization of $f$.

In order to force the monicity condition for Theorem~\ref{cohnthm} we
can do a random shift of the variables $x_i$ by the replacement
$x_i\to x_i+\alpha_i$ for $\alpha_i\in\F$ all picked independently at
random. However, this yields the monicity of the matrix $L$ iff
$f(\alpha_1,\ldots,\alpha_n)\ne 0$ which will fail to hold for a
noncommutative polynomial that vanishes for all scalar substitutions
(like $xy-yx$, for instance). To get around this problem we will
apply Amitsur's theorem (which is a strengthening of the Amitsur-Levitski
theorem).

\begin{theorem}[Amitsur's theorem rephrased]
  Suppose $f\in\FX$ is a nonzero degree $d$ polynomial. Then for
  random $d\times d$ scalar matrices $M_i$ the resulting matrix value
  $f(M_1,M_2,\ldots,M_n)$ is an invertible $d\times d$ matrix with
  high probability. 
\end{theorem}

Thus, we can replace $x_i$ by $X_i+M_i$, where $X_i$ is a generic
$d\times d$ matrix of \emph{noncommuting variables}. The resulting
polynomial $f(X_1+M_1,X_2+M_2,\ldots,X_n+M_)$ yields the desired
monic property in the blown-up linear matrix $L'$ (obtained by applying
the same subsitution in $L$).
